\section{Rappels sur le poulet de chair : appareil digestif et microbiote}

\subsection{Particularités anatomo-physiologiques du tube digestif}

Le poulet de chair présente une organisation digestive adaptée à une ingestion fréquente et à une croissance rapide. La digestion repose sur la complémentarité entre des compartiments de stockage, de sécrétion, de broyage mécanique, de digestion enzymatique et d'absorption. Cette organisation doit être rappelée avant d'interpréter l'effet d'un additif distribué dans l'aliment ou dans l'eau de boisson, car le site de libération, le temps de séjour et l'environnement physicochimique conditionnent l'exposition réelle de la muqueuse et du microbiote \citep{Svihus2014,RavindranAbdollahi2021}.

\subsubsection{Bec, cavité buccale, langue, œsophage et jabot}

Le bec assure la préhension de l'aliment et de l'eau. Chez l'oiseau, la mastication au sens mammalien est absente : la fragmentation est limitée dans la cavité buccale et le bol alimentaire gagne rapidement l'œsophage. Le jabot constitue une dilatation œsophagienne capable de stocker temporairement l'ingéré. Son rôle est particulièrement important lorsque la prise alimentaire est intermittente ou lorsque la vitesse d'ingestion excède momentanément la capacité des segments digestifs distaux. Il participe ainsi à la régularisation du flux alimentaire plutôt qu'à une digestion enzymatique majeure \citep{Svihus2014}.

Pour une substance administrée dans l'eau, cette première portion du tube digestif rappelle que la cinétique d'exposition diffère de celle d'un composé incorporé dans une matrice alimentaire. Le volume d'eau consommé, sa distribution dans la journée et la stabilité du composé en solution deviennent alors des déterminants de l'exposition.

\subsubsection{Proventricule et gésier}

L'estomac aviaire est fonctionnellement divisé entre le proventricule, principalement glandulaire, et le gésier, principalement musculaire. Le proventricule contribue aux sécrétions acides et enzymatiques, tandis que le gésier exerce un rôle mécanique majeur de broyage et de régulation du passage vers l'intestin. La structure physique de l'aliment peut modifier le fonctionnement du gésier et, indirectement, les conditions de digestion en aval \citep{Svihus2014,RavindranAbdollahi2021}.

Cette distinction est importante dans la présente thèse : une poudre de feuilles incorporée à l'aliment peut agir à la fois par sa composition chimique et par sa contribution à la matrice physique, tandis qu'un extrait dissous dans l'eau ne reproduit pas ce contexte. Les résultats obtenus avec ces deux formes ne doivent donc pas être assimilés sans précaution.

\subsubsection{Intestin grêle, cæca et cloaque}

Le duodénum, le jéjunum et l'iléon assurent l'essentiel de la digestion enzymatique et de l'absorption des nutriments. Le pancréas et le foie apportent des sécrétions indispensables à ces processus. Les deux cæca représentent des compartiments particulièrement importants pour les fermentations microbiennes, tandis que le cloaque constitue la portion terminale commune aux voies digestive, urinaire et génitale \citep{RavindranAbdollahi2021,Oakley2014}.


\subsection{Glandes annexes : pancréas et foie}

Le pancréas exocrine fournit des enzymes digestives et des bicarbonates au duodénum, tandis que le foie intervient dans la production biliaire, le métabolisme des nutriments et de nombreuses fonctions de détoxication et de stockage. Une variation d'un métabolite sanguin ou d'un marqueur hépatique après supplémentation ne peut donc être interprétée isolément comme un bénéfice global : elle doit être replacée dans le métabolisme de l'oiseau et dans les autres réponses mesurées.

\subsection{Microbiote digestif du poulet}

Le tube digestif héberge des communautés microbiennes dont la composition varie selon l'âge et le segment anatomique. Les cæca sont particulièrement riches et diversifiés par rapport aux segments proximaux. Les approches modernes de séquençage montrent que le microbiote ne se réduit pas à une opposition entre une flore dite bénéfique et une flore dite nuisible : il s'agit de communautés complexes dont les fonctions peuvent être redondantes et dépendantes du contexte \citep{Oakley2014}.

\subsubsection{Colonisation et évolution avec l'âge}

La colonisation débute très tôt après l'éclosion et évolue rapidement avec l'environnement, l'alimentation, l'âge et les conditions d'élevage. Cette dynamique impose de comparer des prélèvements réalisés au même site et au même âge. Un effet observé dans le contenu cæcal ne peut pas être transposé automatiquement à l'iléon, et une différence observée à un âge donné ne prouve pas une modification persistante.

Les communautés digestives peuvent être étudiées par culture de groupes ciblés ou par méthodes moléculaires donnant une vision plus large de la composition. Ces approches ne fournissent pas la même information. Les différences méthodologiques et leurs conséquences pour les essais sur les feuilles d'olivier sont examinées dans le chapitre consacré aux réponses intestinales \citep{Oakley2014}.
